\subsection{Selección de la biblioteca de compresión}

A partir de la matriz de la tabla \ref{tab:evaluacion_bibliotecas}, se seleccionó
\textbf{Blosc2} como biblioteca de compresión sin pérdida para este trabajo.

La elección se sustenta en que Blosc2 ofrece el entorno más adecuado para materializar
la arquitectura propuesta de almacenamiento comprimido por bloques. Su diseño está
alineado con datos binarios y numéricos en memoria, incorpora filtros reversibles útiles
para arreglos de punto flotante y permite combinar compresión y organización por
fragmentos dentro de un mismo marco de trabajo \cite{blosc2_docs}.

Además, Blosc2 permite conservar flexibilidad experimental sin cambiar de biblioteca.
Dentro de su marco pueden explorarse distintos filtros, configuraciones y códecs, lo que
resulta especialmente valioso en una investigación cuyo interés no es solo comprimir, sino
estudiar el compromiso entre ahorro de memoria, sobrecarga temporal y exactitud
numérica bajo una arquitectura controlada. En consecuencia, Blosc2 se adopta como la
opción más consistente con los objetivos técnicos y metodológicos de este trabajo.
